% ── Postup výzkumu — metodologie a struktura práce ───────────────────────────
% NB: Tato kapitola popisuje vlastní výzkumný postup. Datové zdroje jsou citovány
% v popiscích příslušných obrázků a tabulek; zdroje právních a věcných tvrzení
% jsou citovány v kapitolách, kde jsou použity argumentačně.

Předložená práce analyzuje, zda nízké pokrytí \acins{KS}, nízká odborová organizovanost a~slabé zapojení sociálních partnerů do~aktivních opatření na trhu práce jsou spojeny se zranitelností vůči mzdovým, technologickým a~demografickým tlakům.

\paragraph{Hypotéza práce}\label{par:hypoteza} zní: \emph{Posílení sociálního dialogu a~kolektivního vyjednávání v~\acloc{ČR} by přispělo k~udržitelnému hospodářskému rozvoji.} Udržitelný hospodářský rozvoj je zde chápán ve smyslu strategických cílů vymezených \acs{OECD} a~Agendou~2030 \acs{OSN} -- tedy jako rozvoj doprovázený důstojnou prací, spravedlivým mzdovým přenosem produktivity a~potlačením příjmové nerovnosti. Hypotéza vychází z~předpokladu, že slabý sociální dialog ponechává náklady strukturálních změn trhu práce na~jednotlivých pracovnících, což zvyšuje riziko chudoby, oslabuje mzdovou konvergenci a~omezuje schopnost ekonomiky absorbovat transformace spojené s~automatizací a~stárnutím populace. Reforma v~podobě posílení kolektivního vyjednávání naopak tyto náklady mzdové adaptace mírní a~kolektivně vyvažuje.

\paragraph{Výzkumné otázky}\label{par:vyzkumne_otazky} rozkládají hypotézu do~pěti dílčích vědeckých otázek, na~něž odpovídají jednotlivé části praktické analýzy:
\begin{enumerate}
  \item Jaký je kvantitativní a~právní stav kolektivního vyjednávání v~\acloc{ČR} z~hlediska pokrytí \acpgen{KS}, odborové organizovanosti a~reálné praxe administrativního rozšiřování závaznosti?
  \item Jakým způsobem se smluvní mzdové a~pracovní standardy v~\acloc{ČR} odrážejí v~aspektech sociálního postavení zaměstnanců?
  \item S~jakými hlavními strukturálními nerovnováhami se český trh práce potýká a~jaké jsou jejich souvislosti s~institucionální silou sociálního dialogu?
  \item Jak se v~evropském srovnání profiluje rovnováha zájmů tří hlavních představitelů sociálního dialogu (odborů, zaměstnavatelů a~státu) vyšetřená metodou ternární syntézy?
  \item Jakým způsobem může sociální dialog zprostředkovat adaptaci na technologické, demografické a~odvětvové transformace spojené s~Průmyslem 4.0 a~\acs{AI}?
\end{enumerate}

\paragraph{Cíle práce}\label{par:cile} spočívají v~posouzení stavu a~účinnosti současného sociálního dialogu v~\acloc{ČR} v~evropském srovnání a~ve~vyhledání právních a~organizačních nástrojů přenositelných z~vybraných evropských systémů. Dílčí cíle jsou formulovány tak, aby na sebe logicky navazovaly:
\begin{enumerate}[label = (\roman*),noitemsep]
  \item vymezit právní a~organizační uspořádání sociálního dialogu v~\acloc{ČR} a~stanovit jeho kvantitativní stav;
  \item zasadit český model do~strukturovaného srovnání s~Dánskem, Rakouskem, Německem, Polskem a~Slovenskem;
  \item vyšetřit statistické souvislosti mezi pokrytím \acpgen{KS} a~klíčovými mzdovými, výkonnostními a~časovými ukazateli trhu práce v~panelu zemí \aclgen{EU};
  \item modelovat hypotetické dopady zvýšeného pokrytí \acpgen{KS} v~\acloc{ČR} na národní mzdové a~pracovní parametry za použití regresních odhadů;
  \item konfrontovat makroekonomická zjištění s~empirickým obsahem konkrétních podnikových a~odvětvových kolektivních smluv v~rámci čtyř vybraných případových studií.
\end{enumerate}

\paragraph{Použité metody}\label{par:metody} kombinují deskriptivní statistiku, komparativní analýzu, korelační modelování, právní interpretaci fungování sociálních institucí a~metodu případových studií. Deskriptivní část sleduje časové řady a~mezistátní rozdíly v~hlavních mzdových a~institucionálních ukazatelích. Komparativní srovnání vybraných národních systémů je strukturováno podle tří dimenzí: časové, geografické a~systémové. Statistická část vyšetřuje, zda se odlišná míra pokrytí \acpgen{KS} v~různých evropských zemích pojí s~odpovídajícími mzdovými, časovými a~výkonnostními parametry. Regresní model pro \acacc{ČR} následně slouží ke kvantitativní interpretaci hypotetických scénářů posílení smluvního pokrytí. Závěrečná mikroekonomická rovina, reprezentovaná případovými studiemi, tyto vztahy konfrontuje s~realitou konkrétních PKS a~KSVS.

\paragraph{Výběr komparativních zemí}\label{par:vyber_zemi} je proveden účelově za účelem vyhodnocení klíčových kontrastů. Dánsko představuje severský sociální model (flexicurity), v~němž je vysoká vnější numerická flexibilita trhu práce vyvažována nadstandardní příjmovou jistotou, masivními výdaji na~\aclacc{APZ} a~silnou, téměř plošnou mzdovou regulací prostřednictvím odvětvových smluv. Rakousko je zvoleno jako nejbližší institucionální vzor pro případnou diskuzi o~reformě českého rozšiřování závaznosti, neboť doložitelně kombinuje střední odborovou hustotu s~téměř univerzálním smluvním pokrytím těžícím z~kammer-systému a~odvětvových standardů. Německo reprezentuje klíčového ekonomického partnera \acgen{ČR} a~zároveň průmyslový model s~paralelní tradicí spolurozhodování (\emph{Mitbestimmung}), jehož novodobý vývoj však důrazně varuje před mzdovou segmentací při deregulaci okrajových forem zaměstnání bez tarifní vazby. Polsko a~Slovensko uzavírají srovnání v~roli negativní komparativní skupiny sdílející s~\acloc{ČR} postkomunistické dědictví, decentralizované podnikové vyjednávání a~slabé využití odvětvových standardů.

\paragraph{Časový rámec výzkumu}\label{par:casovy_ramec} se soustředí na období od~vstupu \acs{ČR} do \acs{EU} v~roce 2004 do~současnosti, přičemž klíčové trendy a~změny jsou analyzovány v~kontextu ekonomických cyklů, technologických posunů a~legislativních reforem. Pro komparativní analýzu jsou použity nejnovější dostupné údaje. Část datových sad je dostupná až od pozdějšího data, pak jsou používány od doby, kdy jsou v~nich konzistentně dostupné panelové země.\par
V~tomto časovém rámci je důležitý kontext 3~událostí globálního významu, které měly podstatný vliv na české hospodářství a~trh práce:
\begin{itemize}
  \item velká recese v~roce 2008 zapříčiněná hypoteční krizí v~USA, jejíž následky odezněly v~\acs{ČR} později než v~jiných evropských zemích;
  \item koronavirová recese v~letech 2020--2021 zapříčiněná bojem s~globální pandemií onemocnění Covid-19, která vyvolala vlnu vládních zásahů do~trhu práce a~dočasného omezení ekonomické aktivity, \acs{ČR} však v~této době uchránila zaměstnanost a~vyšla z~krize s~posílenou měnou;
  \item německá recese trvající od roku 2022 zapříčiněná nárustem cen energií v~důsledku propuknutí Rusko-ukrajinské války, v~\acs{ČR} navíc spojená se silnou podporou napadené Ukrajiny a~jejích emigrantů.
\end{itemize}

\paragraph{Struktura práce}\label{par:struktura} postupuje od~pojmového a~teoretického ukotvení k~empirické analýze a~syntéze. První část vymezuje historické souvislosti, pojmovou základnu a~aktuální stav poznání (kap.~\ref{ch:ramec}). Praktický oddíl nejprve představuje zdroje dat a~způsob normalizace klíčových mzdových a~strukturálních veličin (podkap.~\ref{sec:datova_analyza}), hodnotí celkový mzdový a~institucionální stav českého trhu práce (kap.~\ref{ch:stav}), provádí srovnání evropských systémů (kap.~\ref{sec:modely}), vyšetřuje statistické vztahy (kap.~\ref{sec:korelace}) a~výsledná zjištění konfrontuje v~případových studiích (kap.~\ref{sec:case_studies}). Závěrečná syntéza a~hodnocení hypotézy jsou formulovány v~kapitole~\ref{sec:zhodnoceni}.

\paragraph{Limity výzkumu}\label{par:limity} pramení především z~metodologických omezení spojených s~nesnadnou detekcí jednoznačných kauzálních vazeb. Panelová šetření a~harmonizované evropské databáze sice poskytují spolehlivou empirickou základnu, řada ukazatelů však v~rámci srovnání kombinuje administrativní hlášení, výběrová šetření a~modelové extrapolace. Korelační modelování z~povahy věci nemůže izolovat dopad kolektivního vyjednávání od širších vlivů fiskální politiky, odvětvové struktury či historické úrovně společenského rozvoje a~důvěry. Modelové propočty pro \acacc{ČR} pak nenesou predikční váhu pro jednotlivé domácnosti, nýbrž slouží jako srovnání systémových parametrů při standardizované mzdové základně.
